\section{Contexto general}
\label{sec:contexto_general}

El municipio de San Julián, ubicado en la provincia Ñuflo de Chávez del departamento de Santa Cruz, es una zona con intensa actividad productiva, agrícola y comercial. En este municipio, el MÓDULO TECNOLÓGICO PRODUCTIVO SAN JULIÁN cumple un papel muy importante, porque concentra la formación del Bachillerato Técnico Humanístico (BTH) para los jóvenes de secundaria que provienen de las distintas unidades educativas del distrito, siguiendo lo que propone el Modelo Educativo Sociocomunitario Productivo.

Para responder a las necesidades de la región, el Módulo ofrece siete especialidades: Construcción Civil, Contaduría General, Electrónica, Gastronomía, Mecánica Automotriz, Sistemas Informáticos y Textiles y Confecciones. Cada inicio de gestión, esta oferta atrae a cerca de 400 estudiantes de secundaria que desean obtener el título de Técnico Medio. Son tantos postulantes que los espacios en aulas y talleres no alcanzan y la demanda supera cada año la capacidad del Módulo. Esta situación fue confirmada en el trabajo de campo del primer bimestre de 2026 con 60 estudiantes y 5 funcionarios, como establece la normativa educativa \parencite{ministerioeducacion2023}.

Entre 2022 y 2024, la inscripción se realizaba de forma presencial, en ventanillas y con registros escritos a mano. Como los cupos por carrera son limitados y todos querían asegurar un lugar en la especialidad que les gusta, las familias llegaban con días de anticipación para hacer fila a la intemperie, a veces bajo sol o lluvia, con cansancio, pérdida de días de trabajo y mucho desorden.
